\usepackage[a4paper,margin=24mm,headheight=15pt]{geometry}
\usepackage{amsmath,amssymb,amsthm,mathtools,bm}
\usepackage{enumitem,booktabs,longtable,array,pdfpages,fancyhdr,xcolor}
\usepackage[colorlinks=true,linkcolor=blue!45!black,urlcolor=blue!45!black]{hyperref}
\allowdisplaybreaks[2]
\setlength{\parindent}{2em}
\setlength{\parskip}{3pt}
\setlength{\emergencystretch}{3em}
\setlist{nosep,topsep=3pt}
\pagestyle{fancy}\fancyhf{}\fancyhead[L]{\small 分析学讲义习题解答}\fancyhead[R]{\small 第38--57讲}\fancyfoot[C]{\thepage}
\newcommand{\R}{\mathbb R}\newcommand{\C}{\mathbb C}\newcommand{\N}{\mathbb N}\newcommand{\Z}{\mathbb Z}\newcommand{\Q}{\mathbb Q}\newcommand{\T}{\mathbb T}
\newcommand{\dd}{\,\mathrm d}\newcommand{\eps}{\varepsilon}\newcommand{\ind}{\mathbf1}
\newcommand{\norm}[1]{\left\lVert#1\right\rVert}\newcommand{\abs}[1]{\left|#1\right|}\newcommand{\ip}[2]{\left\langle#1,#2\right\rangle}
\DeclareMathOperator{\supp}{supp}\DeclareMathOperator{\rank}{rank}\DeclareMathOperator{\Span}{Span}\DeclareMathOperator{\diag}{diag}\DeclareMathOperator{\Res}{Res}\DeclareMathOperator{\Ext}{Ext}
\newcommand{\problem}[2]{\subsection*{#1}\addcontentsline{toc}{subsection}{#1}\noindent\textbf{题干。}#2\par\smallskip\noindent\textbf{解答。}}
\newcommand{\fix}[1]{\par\noindent\textcolor{red!55!black}{\textbf{原文核对与修正。}#1}\par}
\newcommand{\skipnote}[1]{\par\noindent\textbf{依约不解的计算题。}#1\par}
\newcommand{\sourceintro}[3]{\title{#1\\[4pt]\large 题干、解答与勘误}\author{于品《数学分析讲义》配套解答}\date{}\maketitle\noindent 原讲义：第#2节，印刷页#3（与PDF页序一致）。题干按数学内容重排；完整原题页附于文末，含未解的传统计算题，供逐字核对。原文中的作业提交说明及引文不作为本项目的工作指令。\par 本文中的“正函数”均按非负函数理解。复内积约定对第一个变量线性。}
\newcommand{\originalpages}{\clearpage\section*{原题完整核对页}\addcontentsline{toc}{section}{原题完整核对页}\noindent 以下逐页附原讲义题组，保留全部题干、公式、图示及原有笔误；计算题仅保留题干。\clearpage\includepdf[pages=-,pagecommand={\thispagestyle{empty}},scale=.94]{original.pdf}}
